\section{Results and Analysis}
\label{sec:results}

\subsection{Quantitative Results}
\label{sec:results-quant}

\begin{table}[t]
  \centering
  \caption{Test-set results against the unmodified baseline. \TODO{Mean $\pm$ std over seeds.}}
  \label{tab:main}
  \begin{tabular}{lcc}
    \toprule
    Method & Metric 1 & Metric 2 \\
    \midrule
    Baseline (unmodified) & -- & -- \\
    Ours & -- & -- \\
    \bottomrule
  \end{tabular}
\end{table}

\TODO{Interpret \cref{tab:main}: where the modification helps, where it does not, and why.}

\subsection{Ablation Study}
\label{sec:results-ablation}

\begin{table}[t]
  \centering
  \caption{Ablation: baseline component vs.\ substituted component.}
  \label{tab:ablation}
  \begin{tabular}{lcc}
    \toprule
    Variant & Metric 1 & Metric 2 \\
    \midrule
    Baseline component & -- & -- \\
    Substituted component & -- & -- \\
    \bottomrule
  \end{tabular}
\end{table}

\TODO{Interpret \cref{tab:ablation}, do not only report it.}

\subsection{Empirical Sensitivity Map}
\label{sec:results-sensitivity}

\begin{figure}[t]
  \centering
  \placeholderfig{2.5in}
  \caption{Empirical Sensitivity Map. \TODO{Name the critical variable, the swept values, and the number of runs per value.}}
  \label{fig:sensitivity}
\end{figure}

\TODO{The critical variable, why it was chosen, and what \cref{fig:sensitivity} reveals about performance and stability.}

\subsection{Qualitative Results}
\label{sec:results-qual}

\begin{figure}[t]
  \centering
  \begin{subfigure}{0.48\linewidth}
    \placeholderfig{1.8in}
    \caption{Success cases.}
  \end{subfigure}\hfill
  \begin{subfigure}{0.48\linewidth}
    \placeholderfig{1.8in}
    \caption{Failure cases.}
  \end{subfigure}
  \caption{Qualitative examples. \TODO{Describe what each example shows.}}
  \label{fig:qual}
\end{figure}

\TODO{Explain the success and failure cases in \cref{fig:qual} and what the failures say about the method.}
